% Generated by tools/materialize_round8_referee.py; do not hand-edit.
\section{Round-eight reconstruction: Schur-completed fluctuation geometry and direct cumulant covariance}
\label{sec:r8-b3}

The raw second derivative of the perspective entropy along the nonlinear
manifold $A=A_f$ is not intrinsically positive.  We do not call it a metric.
The covariance is first constructed as the positive limit of exact
finite-volume cumulants and is then identified with the Schur-completed KKT
operator of the constrained variational problem.

\subsection{Raw curvature and the missing constraint multiplier}

For
\[
 \Phi(g,a)=g\log(g/a)-g+a,
 \qquad g=qa,
\]
the chain rule gives
\[
 D^2\Phi=\frac{|\dot g-q\dot a|^2}{qa}
 +(1-q)D^2a.
\]
The last term may have either sign.  Introduce an independent reference
variable $a$ and the nonlinear constraint
\[
 \mathcal C(f,a)=a-A_f=0
\]
with multiplier $\lambda$.  At a stationary phase the bordered Lagrangian is
\[
 \mathscr L=I_0(f_0)+\int\Phi(\Gamma,a)
 +\langle\lambda,a-A_f\rangle+	ext{balance multipliers}.
\]

\begin{lemma}[KKT curvature cancellation]
\label{lem:r8-b3-kkt}
On the critical tangent space, the second derivative of the multiplier term
contributes $-\langle\lambda,D^2A_f[\dot f,\dot f]\rangle$.  The true reduced
Hessian is the Schur complement of the bordered KKT Hessian, not the raw
perspective expression alone.
\end{lemma}

\begin{proof}
Differentiate the constraint twice along a feasible curve:
$\ddot a=DA_f\ddot f+D^2A_f[\dot f,\dot f]$.  In the second variation of the
Lagrangian the terms involving $\ddot f,\ddot a$ vanish by stationarity, while
the displayed curvature is retained with the multiplier sign.  Eliminating
$(\dot a,\dot\lambda)$ by the linearized constraint is precisely a Schur
complement operation.
\end{proof}

\subsection{Covariance from exact finite-volume cumulants}

For a bounded test $F$ let $S_F^\varepsilon$ be the exact centered additive
observable and define
\[
 \Sigma_\varepsilon(F,G)=\mu_\varepsilon^{-1}
 \operatorname{Cov}(S_F^\varepsilon,S_G^\varepsilon).
\]

\begin{theorem}[Positive cumulant limit]
\label{thm:r8-b3-cumulant}
On every regular prepared source chart the forms $\Sigma_\varepsilon$ converge
on a dense nuclear test space to a continuous positive semidefinite form
$\Sigma$.  Its null space consists exactly of collision invariants and dynamic
gauges.
\end{theorem}

\begin{proof}
B2 source sewing gives convergence of the logarithmic moment generating
function in a neighborhood of zero, uniformly with two source derivatives.
Differentiation at the exact finite mean gives convergence of covariances.
Every finite-volume covariance is positive semidefinite, hence so is the
limit.  If $F$ is an invariant or a telescoping dynamic gauge, its centered
sum is a boundary term and has zero scaled variance.  Conversely, zero
variance makes every two-sided pressure derivative vanish; the backward
observability theorem below identifies such a test with an invariant plus a
gauge.
\end{proof}

\subsection{Weighted observability and closed range}

Let $\mathcal B$ be the linearized density--contact balance operator in the
weighted graph Hilbert spaces inherited from B2.  Its adjoint is the backward
kinetic equation with reflected boundary data.

\begin{lemma}[Backward hypocoercive estimate]
\label{lem:r8-b3-observability}
For every finite horizon and regular phase there is $C_T$ such that a backward
solution $\psi$ orthogonal to collision invariants satisfies
\[
 \|\psi\|_{\mathsf H_T}
 \le C_T\bigl(\|\mathcal B^*\psi\|_{\mathsf H_T^*}
 +\|\gamma^-\psi\|_{\mathsf T_T}\bigr).
\]
\end{lemma}

\begin{proof}
Multiply the backward equation by the inverse Maxwellian weight.  Transport is
skew after the Green identity of B2.  The linearized collision form controls
the velocity-orthogonal component by its spectral gap.  Commuting one spatial
derivative with transport and using the collision control gives the missing
macroscopic modes; conservation removes the finite invariant kernel.  Boundary
terms are exactly the trace norm.  Integrate the differential inequality and
apply Gronwall.
\end{proof}

\begin{theorem}[Closed range and exact gauge annihilator]
\label{thm:r8-b3-range}
The balance operator $\mathcal B$ has closed range.  The annihilator of balanced
tangents is the closed span of collision invariants and backward dynamic
gauges.  No further zero-variance classes occur.
\end{theorem}

\begin{proof}
Lemma~\ref{lem:r8-b3-observability} is the closed-range estimate for
$\mathcal B^*$ modulo its finite-dimensional invariant kernel.  The Hilbert
closed-range theorem gives closed range for $\mathcal B$.  Integration by parts
shows that invariants and gauges annihilate every balanced tangent.  The
orthogonal decomposition supplied by the estimate proves the converse.
\end{proof}

\subsection{Schur-completed covariance operator}

Let $\mathsf T$ be the balanced tangent Hilbert space modulo the null space in
Theorem~\ref{thm:r8-b3-range}.  Define the bilinear form
\[
 \mathfrak q^{\rm sc}(u,v)=\Sigma^{-1}(u,v)
\]
initially on the range of the covariance map.

\begin{theorem}[Coercive Schur completion]
\label{thm:r8-b3-schur}
The reduced KKT Schur complement from
Lemma~\ref{lem:r8-b3-kkt} closes to the positive form
$\mathfrak q^{\rm sc}$ on $\mathsf T$.  There are $c,C>0$ such that
\[
 c\|u\|_{\mathsf T}^2\le\mathfrak q^{\rm sc}(u,u)
 \le C\|u\|_{\mathsf T}^2.
\]
The possibly negative raw term $(1-q)D^2A_f$ is not asserted to be positive by
itself.
\end{theorem}

\begin{proof}
At finite volume, constrained Legendre duality identifies the inverse
covariance on the quotient with the Schur complement of the bordered Hessian.
The exact finite covariance is positive.  Uniform source curvature and
Lemma~\ref{lem:r8-b3-observability} give lower and upper bounds independent of
volume.  Mosco convergence of these finite quadratic forms identifies the
limit with $\Sigma^{-1}$ and with the closed KKT Schur complement.
\end{proof}

\begin{corollary}[Lax--Milgram covariance]
\label{cor:r8-b3-lax}
For every continuous cotangent $\xi$ on $\mathsf T$ there is a unique
$u_\xi\in\mathsf T$ such that
\[
 \mathfrak q^{\rm sc}(u_\xi,v)=\langle\xi,v\rangle.
\]
The operator $\xi\mapsto u_\xi$ is the covariance operator $\Sigma$.
\end{corollary}

\begin{proof}
Apply Lax--Milgram to the continuous coercive form in
Theorem~\ref{thm:r8-b3-schur}.  The finite-volume identification and passage to
the limit show that the solution has the cumulant covariance pairings.
\end{proof}

\subsection{Process tightness from connected cumulant densities}

\begin{lemma}[Tree-decaying cumulant density]
\label{lem:r8-b3-tree}
For tests in a bounded nuclear set, the connected $r$-time cumulant density
satisfies
\[
 |\kappa_r(t_1,\ldots,t_r)|
 \le C^r r!\exp\{-c\,\mathsf T(t_1,\ldots,t_r)\},
\]
where $\mathsf T$ is the length of a minimal spanning tree on the time points.
\end{lemma}

\begin{proof}
Expand the source derivatives of the B2 renewal semigroup into connected
collision clusters.  Cutting the largest time gap disconnects a cluster unless
a collision genealogy crosses that gap; the graph semigroup gap gives an
$e^{-c\Delta}$ factor.  Iterating over the edges of a spanning tree gives the
bound.  Contact-time diagonals are integrated as trace measures and are
included in the same estimate.
\end{proof}

\begin{theorem}[Exact-centered Gaussian fluctuation process]
\label{thm:r8-b3-process}
The fields centered at their exact finite-volume means converge in a weighted
nuclear-path topology to the Gaussian process with covariance $\Sigma$.
For adjacent deterministic intervals $I,J$,
\[
 |\operatorname{Cov}(Z(I),Z(J))|
 \le C\int_I\int_J e^{-c|t-s|}\,dt\,ds,
\]
which yields fractional Sobolev tightness and includes near-diagonal contact
correlations.
\end{theorem}

\begin{proof}
Lemma~\ref{lem:r8-b3-tree} makes all cumulants of order at least three vanish
after central scaling and gives the displayed two-interval bound without
assuming a product $|I||J|$ estimate at the common endpoint.  The second
cumulant converges by Theorem~\ref{thm:r8-b3-cumulant}.  Kolmogorov--fractional
Sobolev estimates give scalar tightness; nuclear equicontinuity and Mitoma's
criterion give process tightness.
\end{proof}
